%% [migrated] heading absorbed into chapter wrapper: Introduction: what is lattice QFT?

\begin{itemize}
  \item Fundamentally LQFT is an oxymoron: ``lattice'' refers to a particular class of regulators of a quantum field theory in which the degrees of freedom are restricted to being defined on a discrete set of points, thereby taming divergences in integrals involving them.. However, what we mean by LQFT are the results one extracts after the regulator is removed, that is, the QFT unadorned by the ``lattice'' qualifier
  \item LQFT = QFT defined via a spacetime lattice regulator\aidx{lattice quantum field theory (LQFT)}\aidx{regulator!lattice}
  \item Unlike in condensed matter physics, the lattice here is a mathematical tool with no physical significance
  \item LQFT provides a different though equally valid way to access QFT than your previous classes in DRQFT (= QFT defined via dimensional regularisation\aidx{regulator!dimensional}). The differences make some aspects of QFT more difficult to work with but other aspects easier to consider
  \begin{itemize}
  \item A lattice regulator breaks Poincar\'e symmetries and care is needed to recover it as the regulator is removed
  \item LQFT need not be defined in the context of a perturbative expansion unlike DRQFT for example. Indeed LQFT provides the only known definition of QFT valid for all coupling strengths.
  \item Some QFT concepts are more naturally expressed in LQFT
  \end{itemize}
  \item LQFT is suitable for numerical evaluation of physical quantities in the underlying QFT even at strong coupling:
  \begin{itemize}
  \item LQFT techniques applied to QCD(=LQCD\aidx{lattice QCD (LQCD)}) have a central role in particle and nuclear physics phenomenology and are critical to the determination of the parameters of the Standard Model
  \item LQFT also provides access to non-perturbative information about many other QFTs
  \item Numerically implementing LQFT requires a precise and concrete understanding of what is meant by QFT!!
\end{itemize}
\end{itemize}

These notes are partly based on material covered in the following books
\begin{itemize}
    \item ``Introduction to Quantum Fields on a Lattice,'' (Cambridge Lecture Notes in Physics, 2002) by Jan Smit. 

\item``Quantum Chromodynamics on the Lattice,'' (Springer Lecture Notes in Physics, vol. 788, 2010) by Christof Gattringer and Christian Lang. 

\item ``Quantum Fields on a Lattice,'' (Cambridge Monographs on Mathematical Physics, 1997) by Istvan Montvay and Gernot Munster.  

\item ``Lattice Gauge Theories: An Introduction,'' (World Scientific Lecture Notes in Physics, Vol. 82; 4th Ed. 2012) by Heinz J. Rothe. 
\end{itemize}
